\section{Conclusões}
\label{sec:cap5_conclusoes}

Este capítulo apresenta as conclusões do trabalho, sintetizando os resultados obtidos nas diferentes fases da missão de \gls{rvdb} simulada e discutindo se os objetivos propostos foram atingidos. A partir das análises da transferência orbital, do controle translacional em \gls{lvlh}, do controle de atitude com base flutuante e da atuação do manipulador robótico na aproximação final, são destacados os comportamentos mais relevantes para a operação de \emph{berthing} com perseguidor equipado com manipulador.

Na sequência, são revisitados os objetivos gerais e específicos formulados na introdução da dissertação, explicitando as principais contribuições do modelo e da estratégia de controle adotada, bem como as limitações identificadas. Por fim, são indicadas possibilidades de continuidade do trabalho, com foco em extensões do modelo dinâmico, no aperfeiçoamento das leis de controle e na inclusão explícita da fase de contato e alinhamento fino entre a ferramenta e a porta de acoplamento.

\subsection{Síntese dos objetivos e contribuições}
\label{sec:cap5_sintese_obj_contrib}

O objetivo geral desta pesquisa foi formular e simular computacionalmente a dinâmica de uma espaçonave robótica em órbita, bem como testar algoritmos de controle associados à operação de manipuladores robóticos em ambiente espacial, utilizando MATLAB e Simulink como plataforma principal de desenvolvimento. A partir desse objetivo amplo, o trabalho buscou articular, em um único cenário de missão, a interação entre dinâmica orbital, movimento relativo e atuação de um manipulador montado sobre uma base flutuante.

No que se refere ao primeiro objetivo específico, voltado ao estudo de ferramentas de programação em MATLAB e Simulink, a dissertação desenvolveu um conjunto estruturado de modelos numéricos: dinâmica orbital no referencial ECI, dinâmica relativa em \gls{lvlh}, e dinâmica acoplada base--manipulador descrita por equações de Lagrange. Esses modelos foram integrados em um ambiente de simulação modular, com exportação sistemática de sinais, geração de gráficos e criação de critérios numéricos de desempenho, o que caracteriza um uso aprofundado das ferramentas de modelagem, simulação e pós-processamento disponíveis nesses ambientes.

Em relação ao segundo objetivo específico, que tratava do estudo do ambiente de atividades robóticas (incluindo aspectos de colisão, mobilidade e possíveis danos), o trabalho adotou um cenário simplificado de órbitas circulares coplanares e focou na análise de segurança em termos de janelas de posição e velocidade relativas, limites de atitude, energia cinética e momento angular. Embora não se tenha modelado explicitamente impactos, contato ou falhas estruturais, as restrições impostas às velocidades relativas e aos esforços de controle foram definidas à luz de requisitos de segurança típicos para operações de proximidade, o que permite discutir, em nível de simulação, condições mínimas para reduzir riscos de colisão e manobras abruptas.

O terceiro objetivo específico, associado ao estudo de procedimentos de otimização da mobilidade e à prevenção de falhas operacionais, foi abordado de forma indireta por meio da definição e avaliação de critérios de desempenho. Foram estabelecidas janelas admissíveis de posição e velocidade no referencial \gls{lvlh}, tolerâncias para erros de atitude, limites para energia cinética e momento angular residual, além de uma tolerância para o erro de posição da ferramenta em relação ao ponto desejado na porta de acoplamento. A comparação entre cenários com condições iniciais distintas (cenário A e cenário B) evidenciou a influência de escolhas orbitais iniciais sobre o tempo de missão e o consumo de $\Delta v$, o que se relaciona à ideia de mobilidade mais eficiente e de planejamento de missão menos suscetível a condições desfavoráveis.

Do ponto de vista das contribuições, o trabalho apresentou:
(i) uma modelagem integrada que combina dinâmica orbital no referencial ECI, movimento relativo em \gls{lvlh} e dinâmica base--manipulador em base flutuante;
(ii) uma arquitetura de controle em camadas, englobando transferência orbital, controle translacional no referencial \gls{lvlh}, controle de atitude e controle das juntas do manipulador;
(iii) critérios numéricos explícitos para avaliar posição, velocidade, atitude, energia, momento angular e erro da ferramenta; e
(iv) cenários de simulação que permitem analisar o impacto de diferentes condições iniciais sobre o tempo de missão e o custo propulsivo. Em conjunto, esses elementos atendem ao objetivo geral de estudar, em ambiente de simulação, a operação de manipuladores robóticos em órbita e fornecem uma base coerente para investigações futuras mais detalhadas sobre aproximação e ancoragem com perseguidor em base flutuante.

\subsection{Principais resultados e limitações}
\label{sec:cap5_resultados_limitacoes}

Os resultados numéricos apresentados no Capítulo~\ref{sec:cap4_objetivo_criterios} indicam que a missão é coerente e cumpre os objetivos propostos dentro das hipóteses adotadas.
Na aproximação de longa distância, a transferência de Hohmann conduz o perseguidor a uma órbita próxima com a do alvo, possuindo erro radial e de fase suficientemente pequenos para justificar a descrição posterior do movimento relativo no referencial \gls{lvlh} por meio das equações linearizadas de Hill--Clohessy--Wiltshire.
Essa etapa estabelece condições iniciais consistentes para o início do controle translacional na vizinhança do alvo.

Na aproximação de curta distância, o controle realizado no referencial \gls{lvlh} reduz as componentes de posição e velocidade relativas para dentro de limites preestabelecidos.
Os tempos de estabilização são relativamente longos, mas compatíveis com um compromisso entre rapidez de convergência e limitação de esforços de controle.
A imposição de saturações sobre a velocidade relativa evita que o modelo linearizado produza trajetórias com aproximações excessivamente rápidas, mantendo as velocidades em patamares seguros após a estabilização.

Na aproximação final, a arquitetura de controle mostrou-se capaz de coordenar atitude, movimento da espaçonave e atuação do manipulador.
O controlador de atitude mantém o perseguidor alinhado com o alvo, tanto antes quanto depois da ativação da dinâmica acoplada espaçonave--manipulador, compensando as perturbações geradas pelo movimento das juntas.
Em paralelo, o controle das juntas acompanha a referência gerada pela cinemática inversa e aplicada na dinâmica articular por meio de torques comandados por um controlador \gls{pd}.
Nesta simulação, a referência de posição foi definida como o ponto $(0,0,0)$ no referencial \gls{lvlh}, convertido para o referencial do manipulador robótico (referencial $\{0\}$).
A simulação é encerrada quando a ferramenta atinge essa referência dentro dos limites estabelecidos.

Ressalta-se que o erro de posição da ferramenta não necessariamente converge exatamente a zero, pois o controle empregado é do tipo \gls{pd} (sem ação integral) e a dinâmica numérica inclui saturações e regularizações (por exemplo, limitação de aceleração e regularização da matriz de inércia), além de tolerâncias internas e das aproximações inerentes ao cálculo da cinemática inversa.
Assim, o desempenho é caracterizado pelos valores residuais do erro de posição da ferramenta e por sua estabilização na vizinhança do ponto desejado, conforme os sinais registrados no Simulink.

Ressalta-se, por fim, que as simulações consideram restrições de comando de controle e mecanismos de proteção numérica, mas não visam representar a dinâmica completa de atuação e sensoriamento de uma arquitetura específica de propulsão e de controle de atitude.

\begin{itemize}
    \item A dinâmica relativa é descrita por um modelo linearizado de Hill--Clohessy--Wiltshire em órbitas circulares e coplanares, sem inclusão de perturbações orbitais relevantes (como termos de achatamento $J_2$, arrasto atmosférico ou efeitos de terceiros corpos).

    \item O modelo simulador inclui não-linearidades relevantes por saturação e limitação de sinais (por exemplo, limites de velocidade e aceleração, além de regularizações numéricas nos códigos computacionais).
    No entanto, a cadeia de atuação e sensoriamento não foi fechada com um modelo físico de atuadores e sensores: não há representação explícita de limites de torque/impulso dos atuadores de atitude e translação, nem dinâmica de acionamento, atrasos, ruídos/vieses de medida ou falhas.
    Assim, os resultados devem ser interpretados como desempenho do esquema de controle sob restrições de comando, e não como verificação de viabilidade frente a uma arquitetura específica de atuadores.

    \item O manipulador robótico é modelado sem consideração de flexibilidade estrutural, vibrações elásticas ou efeitos dinâmicos associados a estruturas mais longas ou leves.
\end{itemize}

Essas limitações não invalidam as conclusões obtidas, mas delimitam o escopo de validade dos resultados e motivam estudos futuros em direção a descrições mais completas do ambiente orbital, da interação mecânica na fase de contato e da dinâmica das espaçonaves após o acoplamento.

\subsection{Sugestões para trabalhos futuros}
\label{sec:cap5_trabalhos_futuros}

Os resultados obtidos fornecem base consistente para o estudo de operações de \gls{rvdb} com manipulador robótico em uma espaçonave robótica, mas também indicam possibilidades de extensão do trabalho.
A seguir, apresentam-se sugestões para estudos futuros.

\subsubsection{Extensões do modelo geométrico e dinâmico}

\begin{enumerate}[label=(\alph*)]

    \item \textbf{Modelagem de contato e acoplamento}

    Introduzir modelos de contato entre ferramenta e porta, incluindo rigidez, amortecimento e elementos geométricos, como cones de engate ou guias mecânicas.
    Com isso, seria possível:
    \begin{itemize}
        \item estudar a fase de \emph{soft capture} com controle de força ou impedância no manipulador;
        \item analisar a transferência de impulso e as oscilações induzidas na base durante o acoplamento;
        \item definir tolerâncias posicionais e angulares mais realistas para a interface de engate.
    \end{itemize}

    \item \textbf{Dinâmica orbital com perturbações ambientais e achatamento da Terra}

    Estender o modelo orbital com a inclusão de perturbações, como o termo de achatamento $J_2$, arrasto residual em órbitas baixas e incertezas paramétricas.
    Essa extensão permitiria:
    \begin{itemize}
        \item quantificar a robustez do controlador translacional diante de perturbações sistemáticas;
        \item avaliar desvios acumulados em missões de maior duração;
        \item investigar a necessidade de leis de controle mais sofisticadas para manter as janelas de posição e velocidade.
    \end{itemize}

\end{enumerate}

\subsubsection{Estratégias de controle}

\begin{enumerate}[label=(\alph*),resume]

    \item \textbf{Controle adaptativo ou com ganho variável}

    Investigar controladores com ganhos variando em função da fase da aproximação (por exemplo, distância ao alvo, velocidade relativa, energia cinética ou outros indicadores).
    Essa abordagem permitiria:
    \begin{itemize}
        \item reduzir tempos de estabilização em trechos específicos da missão;
        \item ajustar automaticamente os ganhos de controle conforme o regime dinâmico.
    \end{itemize}

    \item \textbf{Análise formal de estabilidade e robustez}

    Desenvolver análises formais de estabilidade e robustez para o sistema acoplado base--manipulador, por exemplo:
    \begin{itemize}
        \item construção de funções de Lyapunov apropriadas para o sistema completo;
        \item uso de critérios no domínio da frequência para avaliar margens de estabilidade;
        \item derivação de condições de estabilidade diante de variações de massa, inércia, atrasos e perturbações externas.
    \end{itemize}
    Esses estudos forneceriam garantias mais rigorosas para o uso da arquitetura de controle em cenários operacionais mais exigentes.

\end{enumerate}
